Protocol implementations are expected to follow natural-language standards, yet
small semantic deviations in field handling, state-dependent checks, and error
paths can silently undermine interoperability and security. Existing conformance
techniques remain limited by costly manual formalization or weak links between
normative text and code evidence. This paper
presents \tool, a framework that checks protocol implementations against
standards by using protocol fields as the semantic bridge between
natural-language rules and executable code. \tool first searches variables and
extracts field-centric rules from standards, maps each rule to implementation
variables and contexts, and then uses a reasoning-and-test agent to distinguish
consistent behavior, inconsistent behavior, and insufficient evidence. On
real-world implementations of TLS, MQTT, and QUIC, \tool generated
204 conformance-deviation reports. Of these reports, 126 have been confirmed
and fixed by maintainers, roughly nine were false positives, and the remaining
reports are awaiting maintainer response or final adjudication; two confirmed
issues became previously unknown CVEs. Using GPT-5.4 pricing as the cost model,
\tool costs \$1.122 per submitted report, suggesting that field-constrained LLM
reasoning can be practical when paired with code-context extraction and
reproducible validation artifacts.
